\section{Related Work}

\paragraph{Suggestion mining.}
\citet{negi2018open} formalised suggestion mining as the task of identifying
sentences that express a suggestion, defined as a sentence that proposes an
improvement, gives advice, or recommends a course of action. Earlier approaches
relied on hand-crafted rules targeting modal verbs and imperative constructions
\citep{negi2017distant}. The SemEval-2019 Task~9 shared task
\citep{negi2019semeval} established benchmark datasets and attracted 33
participating teams. The best-performing systems used BERT-based architectures,
with the top system achieving 78.1\% suggestion-class F1 on Subtask~A
(in-domain). Subtask~B required cross-domain generalisation from software forums
to hotel reviews, where performance dropped for many systems, indicating that
domain shift is a core challenge in suggestion mining.

\paragraph{Domain adaptation and BERT fine-tuning.}
Pre-trained language models such as BERT \citep{devlin2019bert} learn general
linguistic representations that can be adapted to downstream tasks via
fine-tuning. However, fine-tuned models can be brittle when the target domain
differs from the training domain. \citet{park2019bert} showed that BERT is
unstable for out-of-domain suggestion mining, with high variance across runs on
cross-domain data. \citet{prasanna2019zoho} addressed this by applying
tri-training, a semi-supervised bootstrapping approach, to label target-domain
data without manual annotation, achieving 81.94\% suggestion-class F1 on the
cross-domain subtask. These findings motivate our two-stage fine-tuning
strategy: first training on the larger SemEval dataset for task knowledge, then
adapting to hotel-domain data.

\paragraph{Sentiment analysis versus suggestion mining.}
Suggestion mining is distinct from sentiment analysis. A five-star review is
positive in sentiment but may contain sentences that propose improvements.
Aspect-based sentiment analysis identifies opinions about specific features
(e.g.\ ``the pool was clean''), but does not distinguish between descriptive
praise and constructive suggestions (e.g.\ ``they should heat the pool''). Our
work targets this gap by extracting the subset of sentences that contain
improvement-oriented content, regardless of overall sentiment, and grouping
them thematically using BERTopic \citep{grootendorst2022bertopic} to produce
actionable insights for hotel management.
